% Clase 13 — esfera conductora (1) y conductor puntual (2).
% La transcripcion usa "R" para el radio y para la distancia; aca radio a y
% distancia d, porque el resultado solo depende de d y conviene que se vea.
\begin{tikzpicture}[scale=1]

  % esfera
  \fill[figgray!25] (0,0) circle (1.3);
  \draw[figline, line width=1pt] (0,0) circle (1.3);
  \fill[figgray] (0,0) circle (1.2pt);
  \draw[figvecaux] (0,0) -- (130:1.3);
  \node[figlblaux, fill=figgray!25, inner sep=1pt] at (130:0.72) {$a$};
  \node[figlbl] at (-0.5,-0.55) {$1$};

  % conductor puntual
  \node[figpos] (P) at (5,0) {};
  \node[figlbl] at (5,0.4) {$2$};

  % distancia entre centros
  \draw[figgray, line width=0.6pt, {Latex[length=1.8mm]}-{Latex[length=1.8mm]}] (0,-1.75) -- (5,-1.75);
  \draw[figaux] (0,-1.35) -- (0,-1.9);
  \draw[figaux] (5,-0.15) -- (5,-1.9);
  \node[figlbl, fill=white, inner sep=1.5pt] at (2.5,-1.75) {$d$};

\end{tikzpicture}
